\documentclass[11pt,reqno]{amsart}

%==============================================================================
%  First variations and block selection for Doche's Zhang--Zagier upper bounds.
%  Compile: pdflatex -> bibtex -> pdflatex x2   (uses references.bib).
%==============================================================================

\usepackage[utf8]{inputenc}
\usepackage[T1]{fontenc}
\usepackage{amsmath,amssymb,amsthm,mathtools}
\usepackage[numbers, sort]{natbib}
\usepackage[colorlinks=true,linkcolor=blue,citecolor=blue,urlcolor=blue]{hyperref}
\usepackage{booktabs}
\usepackage[margin=1in]{geometry}
\newcommand{\C}{\mathbb{C}}

\theoremstyle{plain}
\newtheorem{theorem}{Theorem}
\newtheorem{proposition}[theorem]{Proposition}
\newtheorem{lemma}[theorem]{Lemma}
\newtheorem{corollary}[theorem]{Corollary}
\theoremstyle{definition}
\newtheorem{definition}[theorem]{Definition}
\theoremstyle{remark}
\newtheorem{remark}[theorem]{Remark}

\newcommand{\Q}{\mathbb{Q}}
\newcommand{\Z}{\mathbb{Z}}
\newcommand{\R}{\mathbb{R}}
\newcommand{\Qbar}{\overline{\Q}}
\newcommand{\hZ}{h_Z}
\newcommand{\essmin}{\mu^{\mathrm{ess}}(\hZ)}
\newcommand{\rQ}{r_Q}
\newcommand{\dplus}{\partial^{+}}
\newcommand{\dA}{D_A}
\newcommand{\dB}{D_B}
\DeclareMathOperator{\ACT}{ACT}
\DeclareMathOperator{\pp}{pp}
\newcommand{\rb}[1]{r_{#1}}   % normalized active potential of a named block
\newcommand{\repo}{\href{https://github.com/tempcollab/optimizationproblems/tree/main/constants/82a}{\texttt{github.com/tempcollab/optimizationproblems}}}

\begin{document}

\title[First variations for Zhang--Zagier upper bounds]
{First variations and block selection for Doche's Zhang--Zagier upper bounds}

\author{Anonymous Authors}

% \author{Adib Hasan}
% \address{Adib Hasan, SignalPilot AI}

% \author{Akashnil Dutta}
% \address{Akashnil Dutta, Incept Labs}

% \author{Tarik Moon}
% \address{Tarik Moon, SignalPilot AI}

\date{}

\subjclass[2020]{11R06, 11G50, 11Y40 (primary); 49J52, 31A15 (secondary)}
\keywords{Zhang--Zagier height, essential minimum, Mahler measure,
perturbed polynomials, first variations, nonsmooth optimization, interval arithmetic}

\begin{abstract}
Doche's perturbed-polynomial construction gives explicit upper bounds for the
essential minimum $\essmin$ of the Zhang--Zagier height by optimizing real exponent
vectors in a finite dictionary of integer blocks.  Until now, the choice of new blocks
has largely been guided by search.  We recast Doche's value as the nonsmooth ratio
\[
  \log h(q)=\frac{\int_0^1\max(A_q(s),B_q(s))\,ds}{\max(D_A(q),D_B(q))}
\]
and prove a branch-sensitive one-sided first-variation formula for adjoining a new
block $Q$.  The formula is valid in all denominator regimes, including the tie
$D_A=D_B$, where the outer maximum produces a genuine kink.  It yields an explicit
``firing'' criterion: a base-side block lowers the bound to first order exactly when
its active-arc integral is negative, while a perturbing-side block in the
$D_B$-attaining regime fires exactly when its normalized active potential is below the
current value $\log h$.

We then use this criterion in three ways.  First, it predicts two additional
Flammang-table blocks and gives a certified bound
\[
  \essmin\le 0.2538893183,
\]
strictly improving Doche's longstanding $0.25443677$ bound, though not the present
record.  Second, applied diagnostically to the current record of \citet{Gri26}, the
criterion fires on each of its six denominator blocks while natural controls remain
inert.  Third, under an explicit bridge ansatz, it becomes constructive: from the
firing seeds $Q_1,Q_2$ it produces the near-cancellation siblings
$R_0=\pp(Q_1-P_1^5P_2^5P_4P_8)$ and
$R_2=\pp(Q_2+P_1^5P_2^5P_4P_7)$, with firing transferred by a certified
magnitude-subordination inequality.  All inequalities used in certified claims are
verified by outward-rounded interval arithmetic; purely diagnostic grid computations
are labelled as such.  The lower bound remains the $0.2487458$ of \citet{F18}.
Code and certificate references are given at \repo.
\end{abstract}

\maketitle

%==============================================================================
\section{Introduction}\label{sec:intro}
%==============================================================================

For $\alpha\in\Qbar$ let $h(\alpha)$ be the absolute logarithmic Weil height.  The
\emph{Zhang--Zagier height} \citep{Zha95,Zag93} is
$\hZ(\alpha)=h(\alpha)+h(1-\alpha)$, and its essential minimum is
\begin{equation}\label{eq:def}
  \essmin:=\sup\bigl\{\,H\in\R:\ \{\alpha\in\Qbar:\hZ(\alpha)\le H\}\ \text{is finite}\,\bigr\}.
\end{equation}
The value of \eqref{eq:def} is unknown.  The best lower bound currently used here is
$0.2487458=\log 1.282416$ \citep{F18}.  On the upper side, Doche's
$0.25443677=\log 1.289735$ \citep{Doc01b} stood for many years, and the recent certified denominator-first computation of \citet{Gri26} gives
\[
  \essmin\le 0.2536331090.
\]
Thus the presently relevant interval is
\[
  0.2487458\le\essmin\le0.2536331090,
\]
with the lower bound from \citet{F18} and the upper bound from \citet{Gri26}.
The constructions behind the upper bounds are explicit limit-point constructions.
They produce a value $\log h$ from a finite list of integer polynomials and certify
$\essmin\le\log h$.  The central problem is therefore not only to optimize the real
exponents already present in a construction, but also to decide which new integer
polynomials should be adjoined.

This paper develops a local variational theory for that selection problem.  In
Doche's construction one has two branches,
\[
  A(s)=\sum_i q_i\log|P_i(\chi(s))|,\qquad
  B(s)=\log|Q_1Q_2(\chi(s))|+\sum_k q_k\log|Q_k(\chi(s))|,
\]
where $\chi(s)=z(1-z)$ and $z=e^{2\pi i s}$, and the value is the nonsmooth ratio
\begin{equation}\label{eq:ratio}
  \log h(q)=\frac{\Phi(q)}{D(q)},\qquad
  \Phi(q)=\int_0^1\max(A(s),B(s))\,ds,
  \qquad D(q)=\max(\dA(q),\dB(q)).
\end{equation}
Both maxima matter.  The numerator has moving switching points, while the denominator
has a possible one-sided kink when $\dA=\dB$.  The first contribution is to compute the
right derivative of \eqref{eq:ratio} when a block is adjoined at exponent zero.

\begin{theorem}[First variation; proved in Sections~\ref{sec:firstvar}--\ref{sec:unified}]
\label{thm:main}
Let $Q\in\Z[X]$ be adjoined to one branch of an admissible Doche family with exponent
$q_Q\ge0$.  Assume $\log|Q(\chi(\cdot))|\in L^1[0,1]$ and that the two branches do not
coincide on a set of positive measure, the hypotheses \emph{(L)} and \emph{(N)} below.
Then $\log h$ is right-differentiable at $q_Q=0$, and
\begin{equation}\label{eq:main}
  \dplus\!\log h/\partial q_Q\big|_0
  =\frac1D\Bigl[\int_{\ACT}\log|Q(\chi(s))|\,ds-\log h\cdot s_Q\Bigr].
\end{equation}
Here $\ACT$ is the active arc of the branch to which $Q$ is adjoined, and
$s_Q=\deg Q$ if that branch attains $D$ at $q_Q=0$ (at a tie, the perturbed side attains
it for $q_Q>0$) and $s_Q=0$ otherwise.
\end{theorem}

Two points make the proof transparent.  First, the derivative of $\Phi$ is a
dominated-convergence form of the Danskin envelope principle: the moving-boundary term
vanishes because the two branches agree on the switching set.  Second, the term
$-\log h\cdot s_Q$ is contributed solely by the denominator.  This separation is
important at a tie $\dA=\dB$, where the right derivative and left derivative differ by
$\log h\cdot\deg Q/D$.

Theorem~\ref{thm:main} turns block selection into a reduced-cost computation.
Throughout the applications the perturbing branch is the $D$-attaining branch.  If a
new block enters the base branch, it fires exactly when
\begin{equation}\label{eq:intro-base-test}
  \widetilde r_Q:=\int_{\{A>B\}}\log|Q(\chi(s))|\,ds<0.
\end{equation}
If a new block enters the perturbing branch in the same denominator regime, it fires
exactly when
\begin{equation}\label{eq:intro-pert-test}
  r_Q:=\frac1{\deg Q}\int_{\{B>A\}}\log|Q(\chi(s))|\,ds<\log h.
\end{equation}
At an interior optimum this is the complementary-slackness condition
$r_Q=\log h$ for active perturbing blocks, matching the linear-programming duality
framework of \citet{BMQS} restricted to Doche's family.

\begin{corollary}[Selection criterion; proved in Section~\ref{sec:criterion}]
\label{cor:criterion}
For an admissible family in the perturbing-denominator regime, the sign of the
right derivative of $\log h$ after adjoining a new block is given by
\eqref{eq:intro-base-test} for base-side blocks and by \eqref{eq:intro-pert-test} for
perturbing-side blocks.  Thus these two quantities are the correct first-order firing
criteria on the candidate-free anchor.
\end{corollary}

We apply the criterion in three directions.

\emph{A certified improvement over Doche.}  The base-side criterion selects two further
Flammang-table blocks, denoted $j_3$ and $j_9$ below.  After reoptimizing the exponent
vector and certifying the integral by outward-rounded interval arithmetic, we obtain
\begin{theorem}[Certified upper bound; proved in Section~\ref{sec:bound}]
\label{thm:upper}
\[
  \essmin\le 0.2538893183 <0.25443677.
\]
\end{theorem}
This is not the current record; its purpose is to demonstrate that the variational
criterion predicts genuinely improving blocks rather than merely explaining a known
search output.

\emph{A diagnostic for the current record.}  The current record \citep{Gri26} is a
certified denominator-first computation using six denominator blocks, including two
near-cancellation factors $R_0,R_2$ of a degree-$140$ denominator.  The denominator-first
coordinates and Doche's coordinates describe the same surface, since
\[
  \frac{\log M(Q)+\int_0^1(A-B)^+\,ds}{\deg Q}
  =\frac{\int_0^1\max(A,B)\,ds}{\deg Q}.
\]
Consequently the first-variation criterion applies directly to the record.

On the candidate-free anchors for the \citet{Gri26} construction, the regime-correct
one-sided marginal fires on all six denominator blocks and keeps natural control blocks
inert; the values in this diagnostic are fine-grid evaluations, not certified interval
enclosures.  This is stated and proved as Proposition~\ref{thm:record-fires} in
Section~\ref{sec:record}.

\emph{A generator for the near-cancellation factors.}  The diagnostic becomes
constructive once a firing seed and a bridge ansatz are supplied.  Let $Q^*$ be one of
Doche's degree-$28$ distinguished factors and take
$g=P_1^5P_2^5P_4P_{\mathrm{tail}}$, with $P_{\mathrm{tail}}=P_8$ or $P_7$ depending on
the seed.  The primitive part of $Q^*\mp g$ is a new integer block that remains close
to $Q^*$ on the active arc but is coprime to it.

\begin{theorem}[Sibling generator; proved in Section~\ref{sec:generator}]
\label{thm:gen-intro}
For $Q^*\in\{Q_1,Q_2\}$ and the bridge $g=P_1^5P_2^5P_4P_{\mathrm{tail}}$, the
primitive part $R=\pp(Q^*\mp g)$ is a degree-$28$ admissible block, coprime to $Q^*$,
and a certified magnitude-subordination bound transfers firing from $Q^*$ to $R$ on
the candidate-free anchor.  The two resulting blocks are exactly the
near-cancellation factors
\[
  R_0=\pp(Q_1-P_1^5P_2^5P_4P_8),\qquad
  R_2=\pp(Q_2+P_1^5P_2^5P_4P_7)
\]
from the \citet{Gri26} record.  The exponent $5$ is the largest exponent preserving
the monic leading coefficient of the seed.
\end{theorem}

The generator is conditional in a clearly identified way: it explains and certifies the
siblings once the seed and bridge ansatz are supplied; it does not give an unconditional
search-free derivation of the seeds themselves.  Section~\ref{sec:saturation} records
additional finite-search evidence that the cheap base-side gains available to this
particular construction have largely been absorbed.  This evidence is useful for
orientation, but it is not a non-existence theorem.

The rest of the paper is organized as follows.  Section~\ref{sec:functional} fixes the
Doche functional and admissibility conventions.  Sections~\ref{sec:firstvar} and
\ref{sec:unified} prove the first-variation formula, and Section~\ref{sec:criterion}
derives the selection criterion.  Section~\ref{sec:bound} gives the certified
improvement over Doche.  Sections~\ref{sec:record} and \ref{sec:generator} diagnose and
then generate the near-cancellation blocks in the current record.  Section~\ref{sec:verify}
collects the certification protocol, and Appendix~\ref{app:data} lists the polynomial
data needed to reproduce the computations.

%==============================================================================
\section{The Doche functional}\label{sec:functional}
%==============================================================================

If $\alpha$ is an algebraic integer of degree $d$ with minimal polynomial
$P\in\Z[z]$, then
\[
  \hZ(\alpha)=\frac1d\log\bigl(M(P(z))M(P(1-z))\bigr),
\]
where $M$ denotes Mahler measure.  We use the symmetric coordinate
$X=z(1-z)$, invariant under $z\mapsto1-z$.  On the unit circle $z=e^{2\pi i s}$ set
\[
  \chi(s)=z(1-z)=e^{2\pi i s}-e^{4\pi i s},\qquad 0\le s\le1.
\]
The background variational problem is that $\essmin$ can be described as the infimum
of $\int g\,d\mu$, with $g(z)=\log^+|z|+\log^+|1-z|$, over conjugation-invariant
probability measures satisfying the integer-polynomial inequalities
$\int\log|Q|\,d\mu\ge0$ \citep[Thm.~6.5]{BMQS}.  Doche's construction supplies explicit
admissible measures, hence explicit upper bounds, by a perturbed-polynomial limit
process.

Fix Doche's base polynomials $P_1,P_2,P_4,P_6,P_8\in\Z[X]$ and the distinguished
block $Q_1Q_2\in\Z[X]$ of degree $56$ from \citet{Doc01b}.  Additional blocks may be
adjoined either to the base branch or to the perturbing branch.  For nonnegative real
exponents the two branches are
\begin{equation}\label{eq:AB}
  A(s)=\sum_i q_i\log|P_i(\chi(s))|,
  \qquad
  B(s)=\log|Q_1Q_2(\chi(s))|+
       \sum_k q_k\log|Q_k(\chi(s))|,
\end{equation}
and the denominator arguments are
\begin{equation}\label{eq:Dformula}
  \dA=\sum_i q_i\deg P_i,
  \qquad
  \dB=56+\sum_k q_k\deg Q_k .
\end{equation}
Thus the Doche value is the ratio \eqref{eq:ratio}.  A block adjoined to the base side
contributes to $A$ and $\dA$; a block adjoined to the perturbing side contributes to
$B$ and $\dB$.

The use of real exponents is not an extra analytic relaxation.  In Doche's limit
construction, nonnegative real exponent vectors arise as limits of normalized integer
multiplicities in products
$\prod_i P_i^{n_i}\prod_k Q_k^{m_k}$ with $n_i,m_k\in\Z_{\ge0}$, and the value
\eqref{eq:ratio} depends continuously on these proportions.  Hence every
$q\in\R_{\ge0}^{\,\cdot}$ allowed by the construction gives an actual spectrum
limit point and a valid upper bound.

\begin{definition}[Admissible Doche family]\label{def:admissible}
A Doche family is called \emph{admissible} if the hypotheses of
\citet[Lemmas 2--5]{Doc01a} hold.  In the computations below we verify the following
sufficient conditions: the distinguished block has positive degree; each active
perturbing block is squarefree; the active product $W$ of factors with positive
exponent satisfies the required endpoint condition $W(0)=W(1)=1$; and the active
factors are pairwise coprime and coprime to the base, so Doche's nontriviality
condition follows from unique factorization.  Base-side blocks need not satisfy the
endpoint condition; for example $P_1=X$ is a base block.
\end{definition}

\begin{lemma}[Doche's upper-bound lemma, \citealp{Doc01a,Doc01b}]\label{lem:doche}
For every admissible Doche family and every nonnegative real exponent vector in that
family, the value $\log h$ in \eqref{eq:ratio} is a limit point of the
Zhang--Zagier spectrum.  In particular,
\[
  \essmin\le \log h .
\]
\end{lemma}

No optimality of the exponents is needed for Lemma~\ref{lem:doche}.  This is the key
flexibility used below: we may adjoin a candidate block at exponent zero, compute the
right derivative of $\log h$, and then decide whether increasing that exponent can
lower the bound.

%==============================================================================
\section{First variations of the numerator}\label{sec:firstvar}
%==============================================================================

Fix an admissible family with branches $A_0,B$ as in \eqref{eq:AB}, and adjoin a
block $Q\in\Z[X]$ to the \emph{base} branch with exponent $q_Q\ge0$, so
$A(s,q_Q)=A_0(s)+q_Q\log|Q(\chi(s))|$ while $B$ is unchanged. Write
$\Phi(q_Q)=\int_0^1\max\bigl(A_0(s)+q_Q\log|Q(\chi(s))|,\,B(s)\bigr)\,ds$.

Two structural facts are used throughout.  First, $A_0$ and $B$ are finite real
combinations of functions $\log|P(\chi(s))|$; hence they are real-analytic on
$[0,1]$ away from the finite set of contour roots, where they have integrable
$-\infty$ logarithmic singularities.  Second, at such a singular point the vanishing
branch tends to $-\infty$, so a neighbourhood lies in the complementary half
$\{A_0<B\}$ or $\{B<A_0\}$.  These facts hold for every Doche family by construction. 

The lemma then needs only the following, on the single block $Q$ being adjoined and on
the non-degeneracy of the two branches.
\begin{itemize}
\item[(L)] \emph{(Local integrability.)} $\log|Q(\chi(\cdot))|\in L^1[0,1]$. This holds
  whenever $Q$ has no root on the contour $\chi([0,1])$, in which case
  $\log|Q(\chi(\cdot))|$ is bounded; it holds more generally, a contour root
  contributing only an integrable logarithmic singularity. Every block adjoined in this
  paper is contour-root-free.
\item[(N)] \emph{(Non-degeneracy.)} $A_0-B\not\equiv0$ on each analytic sub-arc cut out
  by the structural singularities above; equivalently, the two branches do not coincide
  on a set of positive measure. Since $A_0-B=\sum_i q_i\log|P_i(\chi)|-\sum_k
  q_k\log|Q_k(\chi)|$ is real-analytic on each sub-arc, it either vanishes identically
  there or has only isolated zeros; (N) is the assumption that the former does not
  occur. For \emph{real} exponents this is not a formal consequence of unique
  factorization---it is the statement that the functions $\log|P_i(\chi(\cdot))|$,
  $\log|Q_k(\chi(\cdot))|$ admit no real linear relation on a sub-arc with the given
  coefficients---so we take (N) as a hypothesis, verified for each family used here by
  auditing that the kink set $K=\{A_0=B\}$ is a finite point set: its sign-change count
  is stable across grid refinements, the numerical fingerprint of $A_0-B\not\equiv0$
  (certificate \texttt{firstvar\_02}).
\end{itemize}
We assume throughout this section that the perturbing branch attains the denominator,
$\dA(0)<\dB$, so that $D\equiv\dB$ is locally constant in $q_Q$; this regime
assumption is removed in Section~\ref{sec:unified}, where the formula is proved with no
condition on which branch attains $D$.

\begin{theorem}[First-variation lemma]\label{thm:firstvar}
Under \emph{(L)} and \emph{(N)}, with the perturbing branch attaining $D$, $\Phi$ is
right-differentiable at $q_Q=0$ with
\begin{equation}\label{eq:Phiprime}
  \Phi^{+\prime}(0)=\int_{\{A_0>B\}}\log|Q(\chi(s))|\,ds,
\end{equation}
and consequently $\dplus\!\log h/\partial q_Q|_0=\Phi^{+\prime}(0)/D$.
\end{theorem}

\begin{proof}
For $q_Q\ne0$ set $\psi_q(s)=\bigl[\max(A_0+q_Q\log|Q|,B)-\max(A_0,B)\bigr]/q_Q$,
evaluated at $\chi(s)$.

\emph{Domination.} The map $t\mapsto\max(a+t,b)$ is $1$-Lipschitz, so
$|\psi_q(s)|\le|\log|Q(\chi(s))||$ pointwise; by (L) this dominator lies in $L^1[0,1]$
uniformly in $q_Q$. The bound holds everywhere off the finite structural singular set
$S$ (the contour roots of $A_0,B$, where one branch is $-\infty$); $S$ is null and so
does not affect dominated convergence, and no directional assumption about which branch
is attained there is needed.

\emph{Pointwise limit off the switching set.} Let $K=\{s:A_0(s)=B(s)\}$. For
$s\notin K\cup S$ and $s$ not a contour root of $Q$: if $A_0(s)>B(s)$ then the first
branch dominates for small $|q_Q|$ and $\psi_q(s)\to\log|Q(\chi(s))|$; if $A_0(s)<B(s)$
then $\psi_q(s)\to0$. Thus $\psi_q\to\log|Q(\chi)|\cdot\mathbf 1_{\{A_0>B\}}$ a.e.\ off
$K\cup S$.

\emph{The switching set is null.} On each of the finitely many closed sub-arcs into
which the structural singular set $S$ divides $[0,1]$, $A_0-B$ is real-analytic and, by
(N), not identically zero; a real-analytic function $\not\equiv0$ has isolated
zeros, and isolated zeros in a compact arc are finite in number, so $K$ is finite,
hence null. Thus $K\cup S$ is null.

\emph{Conclusion.} Dominated convergence gives
$\Phi^{+\prime}(0)=\lim_{q_Q\to0^+}\int_0^1\psi_q\,ds
=\int_{\{A_0>B\}}\log|Q(\chi)|\,ds$. The moving boundary of the active set needs no
separate treatment: the a.e.\ limit is already the fixed-set integrand
$\mathbf 1_{\{A_0>B\}}$, and dominated convergence equates the limit of the quotient
with its integral. The first-order boundary term vanishes \emph{because $A_0=B$ on
$K$}---the envelope cancellation---not because the boundary is fixed. Finally
$\log h=\Phi/D$ with $D$ constant near $0$ since the perturbing branch attains $D$.
\end{proof}

The vanishing of the moving-boundary term is the substance of the lemma: a direct
``differentiate under the integral and add a boundary term'' argument must establish
it by hand, whereas here the $1$-Lipschitz bound and the finiteness of $K$ let
dominated convergence absorb it, with the $-\infty$ singularities of $A_0$ controlled
by the same dominator. This is the Danskin/envelope phenomenon for a parametric
integral of a maximum \citep{Ber99}, made rigorous for the present non-smooth
integrand.

%==============================================================================
\section{The unified one-sided marginal}\label{sec:unified}
%==============================================================================

Theorem~\ref{thm:firstvar} assumed the perturbing branch attains $D$. We
remove that regime assumption. Since $\dA,\dB$ are affine (the $D$-formula \eqref{eq:Dformula}),
$D=\max(\dA,\dB)$ is convex and piecewise-linear, with one-sided derivatives
everywhere. At $q_Q=0$ exactly one of three regimes holds:
$\dA(0)>\dB$ (I), $\dA(0)<\dB$ (II), or $\dA(0)=\dB$ (III, a tie).

\begin{theorem}[Unified one-sided marginal]\label{thm:unified}
Let $Q$ be adjoined to side $X\in\{A,B\}$ with exponent $q_Q\ge0$, under
\emph{(L)} and \emph{(N)} and with no condition on which branch attains $D$. Then
$\log h$ is right-differentiable at $q_Q=0$ and
\begin{equation}\label{eq:unified}
  \dplus\!\log h/\partial q_Q\big|_0
  =\frac1D\Bigl[\int_{\ACT(X)}\log|Q(\chi)|\,ds-\log h\cdot s_Q\Bigr],
\end{equation}
where $\ACT(A)=\{A_0>B\}$, $\ACT(B)=\{B>A_0\}$, and $s_Q=\deg Q$ if $X$ attains $D$ at
$q_Q=0$ (at a tie, the perturbed side attains it for $q_Q>0$) and $s_Q=0$ otherwise.
\end{theorem}

\begin{proof}
\emph{The inner derivative is regime-independent.} The proof of
Theorem~\ref{thm:firstvar} uses only the inner $\max(A,B)$: the $1$-Lipschitz
dominator, the a.e.\ limit, and the finiteness of $K=\{A_0=B\}$. It never references
$\dA,\dB$ or which attains $D$. Hence
$\Phi^{+\prime}(0)=\int_{\ACT(X)}\log|Q(\chi)|\,ds$ in every regime (with $\ACT(B)$
for a $B$-side block, the inner branches swapped).

\emph{Quotient rule.} $D(0)\ge56>0$, and both $\Phi,D$ are right-differentiable at
$0$. For $f,g$ right-differentiable with $g(0)\ne0$,
$(f/g)^{+\prime}(0)=\bigl(f^{+\prime}(0)g(0)-f(0)g^{+\prime}(0)\bigr)/g(0)^2$, needing
no two-sided differentiability. With $f=\Phi$, $g=D$, $\Phi(0)/D=\log h$ this is
\eqref{eq:unified} once $s_Q=\dplus D/\partial q_Q|_0$ is identified.

\emph{The denominator's one-sided derivative.} For a max of affine functions the
right derivative is the largest slope over the active index set
\citep[Thm.~23.4]{Roc70}. Only the side $Q$ joined has nonzero slope $\deg Q$; so in
regimes I, II it equals $\deg Q$ if that side attains $D$ and $0$ otherwise. In
regime III both sides attain $D(0)$, and the right derivative is $\max(\deg Q,0)=\deg
Q$: the perturbed side becomes the unique attainer for $q_Q>0$.
\end{proof}

\begin{remark}[The tie is a genuine kink]\label{rem:kink}
In regime III $\log h$ is not two-sided differentiable. For $q_Q\to0^-$ the
unperturbed side attains $D$, so $D$ has left derivative $0$ and $\log h$ has left
derivative $\Phi^{+\prime}(0)/D$ (no cross-term), while the right derivative carries
$-\log h\cdot\deg Q/D$. The two differ by exactly $\log h\cdot\deg Q/D$. Since
admissible perturbations have $q_Q\ge0$, the right derivative is the operative one and
the firing test is its sign. A naive envelope computation that ignores this kink
mispredicts at a tie; \eqref{eq:unified} does not.
\end{remark}

The two cases of \eqref{eq:unified} recover the marginals used in practice: a base
block ($s_Q=0$) gives $\frac1D\int_{\{A_0>B\}}\log|Q|$, the cross-term-free form of
Theorem~\ref{thm:firstvar}; a perturbing block ($s_Q=\deg Q$) gives
$\frac1D\bigl[\int_{\{B>A_0\}}\log|Q|-\log h\cdot\deg Q\bigr]$.

%==============================================================================
\section{The selection criterion}\label{sec:criterion}
%==============================================================================

Throughout this section the perturbing branch attains $D$ (the regime of every family
we construct). The two firing tests of Corollary~\ref{cor:criterion} are the two
cases of \eqref{eq:unified} read off by the side $Q$ joins.

\begin{definition}\label{def:rq}
Let an admissible family have base active set $\{A_0>B\}$ and perturbing active set
$\Omega=\{B>A_0\}$. For a nonzero $Q\in\Z[X]$:
\begin{align*}
  \widetilde r_Q&:=\int_{\{A_0>B\}}\log|Q(\chi(s))|\,ds
  &&\text{(base marginal)},\\
  \rQ&:=\frac1{\deg Q}\int_\Omega\log|Q(\chi(s))|\,ds
  &&\text{(normalized active potential)}.
\end{align*}
\end{definition}

\begin{proof}[Proof of Corollary~\ref{cor:criterion}]
\emph{Base block} ($X=A$): since the perturbing branch attains $D$, a base block
enters the losing argument, so $s_Q=0$ in \eqref{eq:unified} and
\begin{equation}\label{eq:base-marg}
  \dplus\!\log h/\partial q_Q\big|_0=\frac1D\int_{\{A_0>B\}}\log|Q|
  =\frac{\widetilde r_Q}{D},
\end{equation}
with no $\log h$ term; firing is $\widetilde r_Q<0$.

\emph{Perturbing block} ($X=B$): here $Q$ enters the $D$-attaining argument, so
$s_Q=\deg Q$ and $\Phi^{+\prime}(0)=\int_\Omega\log|Q|=(\deg Q)\rQ$, giving
\begin{equation}\label{eq:marginal}
  \dplus\!\log h/\partial q_Q\big|_0=\frac{\deg Q}{D}\bigl(\rQ-\log h\bigr),
\end{equation}
firing iff $\rQ<\log h$. The optimality statements are the Karush--Kuhn--Tucker
conditions for the minimum over $q_Q\ge0$: an interior active perturbing block has
vanishing derivative, hence $\rQ=\log h$; one held at $q_Q=0$ has nonnegative
derivative, hence $\rQ\ge\log h$. In the dual of the \citet[\S6]{BMQS} measure program
restricted to Doche's family, $\rQ-\log h$ is the reduced cost of the block, so these
are its complementary-slackness conditions.
\end{proof}

\begin{remark}[The criterion is a first variation, so the anchor matters]\label{rem:anchor}
Both $\widetilde r_Q$ and $\rQ$ are evaluated at the family \emph{before} $Q$ is
adjoined; the firing test is for an admissible perturbation $q_Q\ge0$ from that anchor.
A family that already contains $Q$ has a different active set and reverses the reading,
and an unconstrained reoptimization driving $q_Q<0$ leaves the admissible cone and
certifies nothing. The sign test is the right derivative on the candidate-free anchor.
\end{remark}

%==============================================================================
\section{A certified improvement over Doche}\label{sec:bound}
%==============================================================================

The base criterion selects two further admissible firing blocks from Flammang's
\citep{F18} table, the low-degree entries of degrees $3$ and $8$ (call them $j_3,j_9$),
each with negative base marginal $\widetilde r_Q<0$ at the held family
$Q_1Q_2\,Q_5^{q_E}Q_6^{q_F}$ (whose perturbing blocks $Q_5,Q_6$ in turn sit on
$\rQ=\log h$). The enriched family is
$h=Q_1Q_2\,Q_5^{q_E}Q_6^{q_F}\,j_3^{q_G}j_9^{q_H}$ with the nine exponents jointly
re-optimized to
\begin{align*}
  q&=(14.011500,13.443930,2.643590,2.299880,0.252420),\\
  q_E&=0.575080,\qquad q_F=0.568800,\qquad
  q_G=0.891590,\qquad q_H=0.066860,
\end{align*}
giving $D=\max(\dA,\dB)=\max(61.65784,72.00176)=72.00176$, the perturbing branch
attaining $D$. Admissibility holds by Lemma~\ref{lem:doche}: $j_9$ is irreducible
hence squarefree, distinct from $j_3$, coprime to $j_3$, and the active dictionary
$\{P_1,P_2,P_4,P_6,P_8,j_3,j_9\}\cup\{Q_1,Q_2,Q_5,Q_6\}$ is pairwise coprime; base
blocks need not satisfy $P(0)=P(1)=1$. Setting $q_H=0$ recovers the previous family
identically, confirming the new block extends a valid bound rather than altering the
integrand.

\begin{proof}[Proof of Theorem~\ref{thm:upper}]
By Lemma~\ref{lem:doche} it remains to enclose $\log h=\Phi/D$ from above. We bound
$\Phi=\int_0^1\max(A,B)\,ds$ by an adaptive quadrature with directed (outward)
rounding, without splitting the maximum: on each cell three guaranteed upper bounds on
$\int_{\text{cell}}\max(A,B)\,ds$ are taken in minimum---a \emph{flat} bound
$(b-a)\max(\sup A,\sup B)$ (never vacuous, since a near-zero of a block sends one
branch to $-\infty$ but the finite branch caps the maximum), an $O(h^2)$
\emph{straddle} bound from $\max(a+x,b+y)\le\max(a,b)+\max(x,y)$ applied to the
midpoint Taylor deviation, and an $O(h^3)$ \emph{midpoint} bound where one branch
dominates---and cells exceeding a deviation cap are bisected, only leaves contributing,
with exact tiling. Every cell bound is rounded toward $+\infty$, and the cells are accumulated by a
\emph{directed} summation: the naive recursive sum of the $n$ nonnegative leaf values is
inflated by the worst-case round-off allowance $\gamma_{n-1}=(n-1)u/(1-(n-1)u)$, $u=2^{-53}$
(the standard error bound for floating-point summation), and rounded up, so the result
provably dominates the exact sum---the accumulated round-off is accounted for, not
absorbed into a single final \texttt{nextafter}. This gives an upper bound
$S\ge\int_0^{2\pi}\max(A,B)\,dt$. The normalizers are divisors of this numerator, so to
keep the quotient an over-estimate they are rounded toward $-\infty$: a value
$\underline{2\pi}\le2\pi$ and a value $\underline D\le D$ obtained by rounding each
nonnegative product $q_k\deg Q_k$ and partial sum of the affine $D$-formula
\eqref{eq:Dformula} downward (the exponents being nonnegative). Then
$\log h=\Phi/D\le S/(\underline{2\pi}\,\underline D)$, with the final division rounded up.
The frontier resolves fully ($742266$ leaves, none unresolved), $S/\underline{2\pi}\le
18.2804777615$ and $\underline D=72.00176$, whence
\[
  \essmin\;\le\;\log h\;\le\;\frac{18.2804777615}{72.00176}=0.2538893183\;<\;0.25443677. \qedhere
\]
\end{proof}

A soundness check compares the per-cell bound against a high-precision reference on
random cells with no violations on the safe side, and a deliberately understated target
is correctly rejected; the independent floating-point value $0.2538891201$ sits just
below the certified enclosure, the correct (outward) direction.

The improvement is modest and is not the record (Table~\ref{tab:gap}); its role is to
certify that the criterion selects blocks predictively, not to settle the constant. The
current record, $0.2536331090$ of \citet{Gri26}, is lower still---and, as the next
section shows, is itself selected block-for-block by the same criterion.

\begin{table}[ht]
\centering\small
\renewcommand{\arraystretch}{1.2}
\begin{tabular}{@{}lll@{}}
\toprule
upper bound & value & gap above $0.2487458$ \citep{F18}\\
\midrule
Doche \citep{Doc01b} & $0.25443677$ & $5.690\times10^{-3}$\\
this paper (Thm.~\ref{thm:upper}) & $0.2538893183$ & $5.144\times10^{-3}$\\
record \citep{Gri26} (Prop.~\ref{thm:record-fires}) & $0.2536331090$ & $4.887\times10^{-3}$\\
\bottomrule
\end{tabular}
\medskip
\caption{Upper bounds against the standing lower bound $0.2487458$ \citep{F18}. This
paper's bound improves Doche's longstanding value; the current record of \citet{Gri26}
is lower still and is diagnosed by the criterion (Section~\ref{sec:record}).}
\label{tab:gap}
\end{table}

%==============================================================================
\section{Diagnostic: the current record is selected by the criterion}\label{sec:record}
%==============================================================================

The criterion is also a diagnostic for constructions that were found by other means. The current record \citep{Gri26},
$\essmin\le0.2536331090$, is reached by a \emph{denominator-first} search: one fixes a
denominator $Q=Q_1Q_2R_0R_2P_7P_9$ of degree $140$ and optimizes the numerator
exponents $q_i$ over $P_1,\dots,P_6,P_8$ by linear programming, where
\begin{equation}\label{eq:R0R2}
  R_0=Q_1-P_1^5P_2^5P_4P_8,\qquad R_2=Q_2+P_1^5P_2^5P_4P_7
\end{equation}
are two \emph{near-cancellation} factors built to sit close to $Q_1,Q_2$. The blocks
$P_7,P_9$ are the degree-$12$ and degree-$16$ Flammang entries that appear as $Q_5,Q_6$
in this paper; $R_0,R_2$ are new and lie outside the dictionary of
Section~\ref{sec:bound}.

\subsection*{One surface, two coordinate systems}
The criterion applies to the record only because the record optimizes the \emph{same}
functional. This identity is the bridge between the two constructions, so we state it.

\begin{proposition}[The record optimizes Doche's functional]\label{prop:samefunctional}
The \citet{Gri26} bound, reported as
$\bigl(\log M(Q)+\int_0^1(A-B)^+\,ds\bigr)/\deg Q$ with $A=\sum_iq_i\log|P_i(\chi)|$ and
$B=\log|Q(\chi)|$, equals Doche's value \eqref{eq:ratio}: by Jensen
$\int_0^1 B\,ds=\log M(Q)$, so
\begin{equation}\label{eq:samefunctional}
  \frac{\log M(Q)+\int_0^1(A-B)^+\,ds}{\deg Q}
  =\frac{\int_0^1\bigl(B+(A-B)^+\bigr)ds}{\deg Q}
  =\frac{\int_0^1\max(A,B)\,ds}{\deg Q}=\frac{\Phi}{D}=\log h.
\end{equation}
\end{proposition}

The denominator-first search descends this surface where the first-variation calculus
differentiates it. Here the perturbing branch is $B$, so $\Omega=\{B>A_0\}$ and
Corollary~\ref{cor:criterion} applies block by block.

\begin{proposition}[First-variation diagnostic for the record]\label{thm:record-fires}
Evaluated on a fine grid (Section~\ref{sec:verify}), the record construction is
diagnosed by the criterion as follows. All six denominator blocks are \emph{perturbing}
(B-side) blocks, so by Theorem~\ref{thm:unified} each marginal is read on the arc
$\{B>A_0\}$. At the candidate-free anchor---the record family with one block removed---the
regime is set by the numerator degree $\dA=125.24$ against the reduced denominator degree
$\dB$: removing a degree-$28$ block or $P_9$ leaves $\dB\le124<\dA$, the \emph{base}
regime, where $s_Q=0$ and the entry test of Theorem~\ref{thm:unified} is
$m_Q:=\int_{\{B>A_0\}}\log|Q(\chi)|\,ds<0$ (no $\log h$ term); only removing $P_7$ leaves
$\dB=128>\dA$, the perturbing regime with test $m_Q-\log h\deg Q<0$. Every one of the six
fires: the active-arc integral $m_Q$ is strictly negative throughout
(Table~\ref{tab:recordfire}), and the two new near-cancellation factors $R_0,R_2$ fire as
strongly as the established blocks. The test discriminates: natural controls drawn from
the base family and from the bounded-coefficient search family have \emph{positive}
active-arc integral and so do not fire---$m_{P_3}=+0.79$, $m_{P_5}=+2.05$,
$m_{X^2-X+1}=+0.54$ (full anchor). At the record's own \emph{full} optimum, where
$\dB=140>\dA$ and the perturbing regime holds, the six active blocks cluster at a common
reduced-cost level $\rQ-\log h\in[+0.0069,+0.0084]$ (spread $0.0015$): not exact KKT
stationarity, but a diagnostic analogue, the multiplicities being fixed at one rather
than freely optimized.
\end{proposition}

\begin{table}[ht]
\centering\small
\renewcommand{\arraystretch}{1.2}
\begin{tabular}{@{}lrrcrrc@{}}
\toprule
block & degree & $\dB$ (cand.-free) & regime & $m_Q$ (cand.-free) & $\rQ-\log h$ at opt. & fires?\\
\midrule
$Q_1$ & $28$ & $112$ & base & $-1.015$ & $+0.0069$ & yes\\
$Q_2$ & $28$ & $112$ & base & $-0.980$ & $+0.0075$ & yes\\
$R_0$ & $28$ & $112$ & base & $-0.996$ & $+0.0073$ & yes\\
$R_2$ & $28$ & $112$ & base & $-0.958$ & $+0.0073$ & yes\\
$P_7=Q_5$ & $12$ & $128$ & pert. & $-0.110$ & $+0.0084$ & yes\\
$P_9=Q_6$ & $16$ & $124$ & base & $-0.373$ & $+0.0077$ & yes\\
\bottomrule
\end{tabular}
\medskip
\caption{The six denominator blocks of the \citet{Gri26} record ($\dA=125.24$,
$\log h=0.25363$). Each is a perturbing (B-side) block, so its marginal is read on the
arc $\{B>A_0\}$ (Theorem~\ref{thm:unified}). \emph{Regime / $m_Q$}: at the candidate-free
anchor, removing the block leaves the reduced denominator degree $\dB$; for $\dB<\dA$
(base regime, $s_Q=0$) the entry test is the active-arc integral
$m_Q=\int_{\{B>A_0\}}\log|Q(\chi)|\,ds<0$, and all six are strictly negative, so each
lowers the bound when adjoined. \emph{At optimum}: at the full optimum $\dB=140>\dA$, the
perturbing regime holds and the reduced cost $\rQ-\log h$ clusters near a common level.}
\label{tab:recordfire}
\end{table}

The diagnostic rests on the exact identity \eqref{eq:samefunctional} and on the unified
marginal of Theorem~\ref{thm:unified}, with the entry test taken in its regime-correct
form: at the candidate-free anchor the perturbing branch no longer attains $D$ for five
of the six blocks ($\dB<\dA$), so the firing condition is $m_Q<0$, the $s_Q=0$ case---not
the normalized $\rQ<\log h$, which is valid only when the perturbing branch attains $D$
(as at the full optimum and the $P_7$-removed anchor). The two readings agree in verdict
here because $m_Q$ is itself strictly negative, but only the $s_Q=0$ form is licensed in
the base regime. The candidate-free anchor is the one called for by
Remark~\ref{rem:anchor}: $m_Q$ is the numerator of the right derivative of $\log h$ as
that block enters at $q_Q=0^+$. The grid computation reproduces the search's
\emph{candidate} value $\log h=0.2536269$ from \eqref{eq:samefunctional} to $10^{-9}$;
this is the uncertified LP value, of which $0.2536331090$ is the slightly weaker
outward-rounded interval enclosure (the certified record of \citet{Gri26}), so the
diagnostic operates on the true optimum and the small gap is the certification margin.

The reading of this section is a \emph{diagnostic}: the criterion fires on $R_0,R_2$
after the fact. The next section strengthens it to a \emph{construction}---given the
distinguished blocks $Q_1,Q_2$ and the bridge ansatz, $R_0,R_2$ are produced, with the
firing transferred to them by a certified inequality, rather than read off a grid.

%==============================================================================
\section{Generating near-cancellation siblings}\label{sec:generator}
%==============================================================================

The diagnostic of Section~\ref{sec:record} fires on $R_0,R_2$ after the fact.
This section explains how the same variational idea becomes constructive once a seed and a bridge ansatz are specified: given a firing perturbing block $Q^*$ and an
admissible \emph{bridge} polynomial $g$, the primitive part $R=\pp(Q^*\mp g)$ is an
explicit degree-preserving coprime firing sibling of $Q^*$, and $R_0,R_2$ are its
instances. The construction is generative in the precise sense that its input is
$Q^*$ and the bridge ansatz---not $R_0,R_2$---and its output is the family they belong
to. What stays heuristic is named in the scope remark: the \emph{choice} of seed and
bridge, not the fact that the recipe yields firing siblings.

Throughout, fix one \emph{anchor family} $F$: the record denominator with \emph{both}
the seed block $Q^*$ and the sibling $R$ removed, with active arc
$\Omega_F=\{B_F>A_F\}$. Removing both is what makes $\Omega_F$ a candidate-free anchor
for each of $Q^*$ and $R$: by Remark~\ref{rem:anchor} the firing test for a block is the
right derivative as that block \emph{enters} at exponent $0$, so the block must be absent
from the family it is scored against. Removing two degree-$28$ blocks leaves the reduced
denominator degree $\dB=84$, well below the numerator degree $\dA=125.24$, so $F$ is in
the \emph{base} regime: by Theorem~\ref{thm:unified} the entry test for a perturbing
($B$-side) block $Q$ is the active-arc integral
\begin{equation}\label{eq:mdef}
  m_Q:=\int_{\Omega_F}\log|Q(\chi(s))|\,ds<0\qquad(s_Q=0),
\end{equation}
the $s_Q=0$ form, not the normalized $\rQ<\log h$ (which would apply only if the
perturbing branch attained $D$). Both $Q^*$ and $R$ are scored on this single
$\Omega_F$, the integrals $m_{Q^*},m_R$ are unambiguous, and ``$Q^*$ fires'' reads
$m_{Q^*}<0$ as a genuine entry test.

\subsection*{Degree preservation and the maximal exponent}
The bridge ansatz is $g_a=P_1^aP_2^aP_4\cdot P_{\mathrm{tail}}$ with
$P_{\mathrm{tail}}\in\{P_7,P_8\}$ of degree $12$. The displacement degree is a clean
count.

\begin{proposition}[Maximal degree-preserving displacement]\label{prop:degfloor}
For $g_a=P_1^aP_2^aP_4\,P_{\mathrm{tail}}$ with $\deg P_{\mathrm{tail}}=12$ one has
$\deg g_a=2a+16$. Hence, since $\deg Q^*=28$,
\[
  2a+16<28\ \Longleftrightarrow\ a\le5\ \Longrightarrow\ \deg\bigl(Q^*\mp g_a\bigr)=28 ,
\]
with the top $28-(2a+16)$ coefficients of $Q^*$ preserved; taking the primitive part does
not change the degree. So $a=5$ is the largest exponent at which the displacement is
guaranteed to fix the degree and leading coefficient \emph{for either sign}, and it is
the maximal exponent preserving the \emph{monic leading coefficient} of $Q^*$ on both
sides; the boundary $a=6$ behaves differently for the two signs, as the proof records.
The record factors $R_0,R_2$ are the instances $a=5$.
\end{proposition}

\begin{proof}
$\deg(P_1^aP_2^aP_4)=2a+4$, so $\deg g_a=2a+16$. For $a\le5$, $\deg g_a<28=\deg Q^*$, so
$Q^*\mp g_a$ has the leading coefficient of $Q^*$ and agrees with it in the top
$28-\deg g_a$ coefficients; the primitive part rescales by a nonzero integer content and
fixes the degree. At $a=6$, $\deg g_6=28$ and the $X^{28}$ coefficient of $Q^*\mp g_6$ is
$\operatorname{lc}(Q^*)\mp\operatorname{lc}(g_6)=1\mp1$: for the subtraction it is $0$ and
the degree drops to $27$, while for the addition it is $2$ and the degree stays $28$ with
a non-monic leading coefficient. In both cases $a=5$ is the last exponent preserving the
monic leading coefficient of $Q^*$.
\end{proof}

For $a=1,\dots,5$ the sibling $R_a=\pp(Q_1-g_a)$ is squarefree, coprime to $Q_1$,
satisfies $R_a(0)=R_a(1)=1$, and has content $1$; these are finite \texttt{sympy} checks
(certificate \texttt{firstvar\_08}), so each $R_a$ is admissible
(Lemma~\ref{lem:doche}). Table~\ref{tab:degfloor} records the family and its degree
boundary.

\begin{table}[ht]
\centering\small
\renewcommand{\arraystretch}{1.2}
\begin{tabular}{@{}rrrrccr@{}}
\toprule
$a$ & $\deg g_a=2a+16$ & $\deg R_a$ & content & admissible & $\deg R_a=28$ & $\rb{R_a}(\Omega_F)$ (per deg.)\\
\midrule
$1$ & $18$ & $28$ & $1$ & yes & yes & $-0.03231$\\
$2$ & $20$ & $28$ & $1$ & yes & yes & $-0.03917$\\
$3$ & $22$ & $28$ & $1$ & yes & yes & $-0.04442$\\
$4$ & $24$ & $28$ & $1$ & yes & yes & $-0.04798$\\
$5$ & $26$ & $28$ & $1$ & yes & yes\quad($R_0=R_5$) & $-0.04982$\\
$6$ & $28$ & $27$ & $1$ & yes & no & $(-0.05233)$\\
\bottomrule
\end{tabular}
\medskip
\caption{The bridge-exponent family $R_a=\pp(Q_1-P_1^aP_2^aP_4\,P_8)$, its degree
boundary (Proposition~\ref{prop:degfloor}), and the per-degree active potential
$\rb{R_a}(\Omega_F)$ (whose ordering, at fixed degree $28$ in the base regime, is that of
the firing quantity $m_{R_a}$). The degree is preserved exactly for $a\le5$; $a=5$ is
Grinsztajn's $R_0$. Over the degree-preserving range $a\le5$ the potential
$\rb{R_a}$ is strictly decreasing (firing strength strictly increasing),
maximised at $a=5$ (Proposition~\ref{prop:restricted-opt}); the strict ordering is
\emph{certified} by the four signed difference enclosures $D_a<0$ of
\texttt{firstvar\_09}. The $a=6$ entry is parenthesised: its margin is the
\emph{unconstrained} optimum but $\deg R_6=27$ breaks degree preservation and the
shared-normalizer identity \eqref{eq:transfer-id}, so it lies outside the comparison
(the margin values for $a\le5$ are reproducible float midpoints; only their strict
ordering is the certified claim). Degree data verified by \texttt{firstvar\_08} with
exact integer arithmetic.}
\label{tab:degfloor}
\end{table}

\subsection*{The marginal-transfer identity and the firing-transfer bound}
In the base regime of $\Omega_F$ the firing quantity is the un-normalized active-arc
integral $m_Q=\int_{\Omega_F}\log|Q(\chi)|\,ds$ of \eqref{eq:mdef}; scoring $Q^*$ and $R$
on the single $\Omega_F$ gives the exact identity
\begin{equation}\label{eq:transfer-id}
  m_R-m_{Q^*}=\int_{\Omega_F}\log\frac{|R(\chi(s))|}{|Q^*(\chi(s))|}\,ds .
\end{equation}
Writing $R=c^{-1}(Q^*\mp g)$ with $c=\operatorname{cont}(Q^*\mp g)\in\Z_{>0}$,
\[
  \log\frac{|R|}{|Q^*|}=-\log c+\log\Bigl|\,1\mp\frac{g}{Q^*}\Bigr| .
\]
Fix $\theta=\tfrac12$ and split $\Omega_F$ into a \emph{bulk}
$\Omega_0=\Omega_F\cap\{|g/Q^*|\le\theta\}$ and a \emph{well}
$W=\Omega_F\cap\{|g/Q^*|>\theta\}$, $\delta:=|W|$.

\begin{lemma}[Magnitude-subordination / firing-transfer bound]\label{lem:transfer}
Let $Q^*$ be a perturbing block with $\mu^*:=\inf_{\Omega_F}|Q^*|>0$, $g\in\Z[X]$ with
$\deg g<\deg Q^*$, and $R=\pp(Q^*\mp g)$ contour-root-free
($\mu_R:=\inf_{\Omega_F}|R|>0$). Then
\begin{equation}\label{eq:transfer-bound}
  \bigl|m_R-m_{Q^*}\bigr|
  \;\le\;\log c\,|\Omega_F|
  +\frac{\theta}{1-\theta}\,|\Omega_0|
  +\int_{W}\bigl|\log|R/Q^*|\bigr|\,ds\;=:\;\mathrm{RHS},
\end{equation}
and the well integral $\int_{W}|\log|R/Q^*||\,ds$ is finite because $\mu_R>0$, $\mu^*>0$,
and $\sup_{\Omega_F}\log|Q^*|<\infty$. Consequently, in the base regime, if
$m_{Q^*}(\Omega_F)+\mathrm{RHS}<0$ then $m_R(\Omega_F)<0$: the sibling $R$ fires.
\end{lemma}

\begin{proof}
Start from \eqref{eq:transfer-id} and bound the integrand in absolute value on each
piece. The content term contributes $-\log c$ uniformly; since $c\in\Z_{>0}$ we have
$\log c\ge0$, so $|{-}\log c|=\log c$ and this term contributes $\log c\,|\Omega_F|$ in
absolute value. On $\Omega_0$, with $u=g/Q^*$ and $|u|\le\theta<1$,
\[
  \bigl|\log|1\mp u|\bigr|\le-\log(1-|u|)\le\frac{|u|}{1-|u|}\le\frac{\theta}{1-\theta},
\]
the first inequality because $1-|u|\le|1\mp u|\le 1+|u|$ and
$\log(1+|u|)\le-\log(1-|u|)$, the second the standard bound
$-\log(1-x)\le x/(1-x)$ for $x\in[0,1)$, the third monotonicity; integrating gives
$\frac{\theta}{1-\theta}|\Omega_0|$.

On $W$ the ratio $|u|=|g/Q^*|$ may exceed $1$, so $|\log|R/Q^*||$ has no useful global
cap and we retain the integral $\int_W|\log|R/Q^*||\,ds$. \emph{Integrability} is
immediate: the integrand $\bigl|\log|R|-\log|Q^*|\bigr|$ is bounded pointwise by
$\log(1/\mu_R)+\log(1/\mu^*)+\sup_{\Omega_F}\log|Q^*|<\infty$---the lower bound
$\mu^*>0$ on $|Q^*|$ controls $-\log|Q^*|$ on the well, where $|Q^*|$ itself dips, just
as $\mu_R>0$ controls $-\log|R|$. This $L^\infty$
bound secures finiteness but is far too crude to certify firing, which is why we bound
the integral cell-by-cell instead. On a contour cell where the directed-rounded
enclosures give $\log|R|\in[\ell_R,L_R]$ and $\log|Q^*|\in[\ell_{Q^*},L_{Q^*}]$, the
integrand's supremum is at most $\max\{L_R-\ell_{Q^*},\,L_{Q^*}-\ell_R\}$, a valid
per-cell bound whether or not $\log|R|>\log|Q^*|$ there. As a cell shrinks around a
shared well of $R,Q^*$ this bound grows only like $\log(1/\text{width})$, so its
width-weighted contribution vanishes; hence the sum of the per-cell caps is a sound
upper bound on $\int_W|\log|R/Q^*||\,ds$ that converges to it as the well cells are
refined, driving the certified integral to the true marginal shift ($\approx10^{-2}$,
cf.\ \eqref{eq:transfer-id}).
Summing the three pieces gives \eqref{eq:transfer-bound}; the firing conclusion combines
it with the base-regime entry test \eqref{eq:mdef} (Theorem~\ref{thm:unified}, $s_Q=0$)
applied to $R$ on $\Omega_F$.
\end{proof}

\begin{theorem}[Sibling generator]\label{thm:generator}
Let $Q^*\in\{Q_1,Q_2\}$ be a degree-$28$ factor of Doche's distinguished block, firing
on the anchor $\Omega_F$, and let $g=P_1^5P_2^5P_4\cdot P_{\mathrm{tail}}$ with
$P_{\mathrm{tail}}=P_8$ (resp.\ $P_7$). Then $R:=\pp(Q^*\mp g)$, with the sign chosen as
in \eqref{eq:R0R2}, is a degree-$28$ admissible block, coprime to $Q^*$, squarefree,
with $R(0)=R(1)=1$, content $1$, and contour-root-free; and the firing-transfer bound
\eqref{eq:transfer-bound} clears the firing gap in the base regime,
\[
  m_R(\Omega_F)\;\le\;m_{Q^*}(\Omega_F)+\mathrm{RHS}\;<\;0 ,
\]
so $R$ fires and is adjoinable as a denominator factor independent of $Q^*$. The
Grinsztajn near-cancellation factors $R_0,R_2$ are exactly the instances
$Q^*\in\{Q_1,Q_2\}$ of this construction.
\end{theorem}

\begin{proof}
Admissibility and degree preservation are Proposition~\ref{prop:degfloor} and the
finite \texttt{sympy} checks (content $1$, coprimality, squarefreeness, $R(0)=R(1)=1$),
certificate \texttt{firstvar\_08}. The contour-root-freeness $\mu_R>0$, the bulk
constants $|\Omega_0|,m_{Q^*}(\Omega_F)$, and the cell-wise well integral
$\int_W|\log|R/Q^*||\,ds$ entering \eqref{eq:transfer-bound} are each a directed-rounded
interval enclosure over $\Omega_F$, computed by the same outward-rounded cell machinery
as the bound of Section~\ref{sec:bound}. The well cells are adaptively bisected on two
criteria---until $|R|^2$ is certified strictly positive (so $\mu_R>0$ is established, not
measured) \emph{and} until each cell's contribution to $\int_W|\log|R/Q^*||\,ds$ falls
below a fixed tolerance---so the certified well integral is the rigorous per-cell bound
of Lemma~\ref{lem:transfer}, with no global cap on $|\log|R/Q^*||$ invoked. Two points
secure the certificate. First, $\Omega_F$ is the candidate-free anchor of
Section~\ref{sec:generator}: \emph{both} $Q^*$ and $R$ are removed from the record
denominator, leaving $\dB=84<\dA$, so the base-regime test $m<0$ of \eqref{eq:mdef}
applies and $m_{Q^*}(\Omega_F),m_R(\Omega_F)$ are genuine entry tests
(Remark~\ref{rem:anchor}), not scores against a family already containing $R$. Second,
the upper bound on $m_{Q^*}=\int_{\Omega_F}\log|Q^*|\,ds$ is interval-sound on cells
whose membership in $\Omega_F$ is uncertain: a cell certainly in $\Omega_F$ contributes
its signed $\int\log|Q^*|$, but a straddle cell---which may lie outside $\Omega_F$, where
its true contribution is $0$---is clamped to its non-negative part, since banking a
negative value would understate (not over-state) the integral. With these
certified constants,
\[
  m_{Q_1}(\Omega_F)\le-1.39,\quad \mathrm{RHS}\le0.10,\quad
  m_{Q_1}+\mathrm{RHS}\le-1.29<0
\]
for $R_0$, and $m_{Q_2}(\Omega_F)\le-1.39$, $\mathrm{RHS}\le0.10$,
$m_{Q_2}+\mathrm{RHS}\le-1.29<0$ for $R_2$. The firing margins $|m_R|\ge1.3$ are
large---the firing is not delicate. Lemma~\ref{lem:transfer} then gives $m_R<0$, and the
base-regime entry test \eqref{eq:mdef} that $R$ fires. Coprimality $\gcd(R,Q^*)=1$ makes
$R$ an independent admissible factor.
\end{proof}

The bound \eqref{eq:transfer-bound} is not tight---its certified value $\approx0.1$
sits well above the true marginal shift $\approx0.02$ from
\eqref{eq:transfer-id}---but with $1.3$ of margin to spare it certifies firing
comfortably. The mechanism is the magnitude-engineering of the bridge: $g$ is small
relative to $Q^*$ on the bulk
$\Omega_0$ ($|g/Q^*|\le\theta$ on measure $\approx0.40$ of $\Omega_F$, and far smaller
than $\theta$ over most of it), so $R\approx Q^*$ there and inherits $Q^*$'s reduced
cost, while in the deep wells $W$ (measure $\delta\approx0.07$) the bridge is comparable
to $Q^*$ and displaces the roots, making $R$ coprime to $Q^*$. The sibling is the same
size as $Q^*$, not an approximation to zero; it fires by inheritance, not by smallness.

\begin{remark}[What is generated, and the open problem]\label{rem:gen-scope}
Theorem~\ref{thm:generator} generates \emph{from a seed}: given $Q^*$ and the bridge
ansatz it produces the firing sibling family containing $R_0,R_2$, every constant a
directed-rounded enclosure rather than a grid reading. It converts the diagnosis of
Section~\ref{sec:record} into the construction ``$R_0,R_2$ are the maximal
degree-preserving coprime siblings of $Q_1,Q_2$, firing by magnitude-subordination of
the bridge.'' What stays heuristic is the choice of seed and bridge; producing the seed
with no input---and showing it optimal among generators---is the open problem.
\end{remark}

\subsection*{The bridge exponent is selected, not merely admitted}
Proposition~\ref{prop:degfloor} shows $a=5$ is the \emph{maximal monic-preserving}
bridge exponent, but admissibility does not by itself single $a=5$ out among the five
candidates $a\in\{1,\dots,5\}$. The next proposition shows the criterion does single it
out: among the monic-preserving (degree-$28$) siblings the firing strength is strictly
increasing in $a$, so $a=5$ is the firing-\emph{optimal} such displacement.

\begin{proposition}[Restricted bridge optimality]
\label{prop:restricted-opt}
Let $R_a=\pp(Q_1-P_1^aP_2^aP_4\,P_8)$ for $a\in\{1,\dots,5\}$ be the degree-preserving
admissible siblings of Table~\ref{tab:degfloor} (each $\deg R_a=28$, content $1$,
squarefree, coprime to $Q_1$, contour-root-free). Then the base-regime firing quantity
$m_{R_a}=\int_{\Omega_F}\log|R_a(\chi)|\,ds$ on the shared anchor $\Omega_F$ is strictly
decreasing in $a$,
\[
  m_{R_1}(\Omega_F)>m_{R_2}(\Omega_F)>m_{R_3}(\Omega_F)
  >m_{R_4}(\Omega_F)>m_{R_5}(\Omega_F),
\]
so the firing strength $|m_{R_a}(\Omega_F)|$ is maximised at the degree boundary
$a=5$. Thus $a=5=$ Grinsztajn's $R_0$ is the firing-optimal exponent among the
degree-preserving admissible siblings.
\end{proposition}

\begin{proof}
Since $\deg R_a=\deg R_{a+1}=28$ and both blocks are scored on the same fixed
$\Omega_F$ in the same base regime, the marginal-transfer identity
\eqref{eq:transfer-id}, applied with $Q^*:=R_a$ and $R:=R_{a+1}$ (note
$\deg(R_{a+1}-R_a)<28$, so Lemma~\ref{lem:transfer} applies verbatim), gives the exact
\emph{consecutive-difference identity}
\begin{equation}\label{eq:diff-id}
  D_a:=m_{R_{a+1}}(\Omega_F)-m_{R_a}(\Omega_F)
     =\int_{\Omega_F}\log\frac{|R_{a+1}(\chi(s))|}{|R_a(\chi(s))|}\,ds,
  \qquad a=1,2,3,4 .
\end{equation}
It therefore suffices to certify $D_a<0$ for $a=1,2,3,4$; the four strict inequalities
chain to the displayed ordering. We certify a rigorous outward-rounded \emph{upper}
bound $D_a^{\mathrm{hi}}$ on each $D_a$ and check $D_a^{\mathrm{hi}}<0$.

The integrand of \eqref{eq:diff-id} is the \emph{signed} difference
$\log|R_{a+1}|-\log|R_a|$; it is not one-signed on $\Omega_F$ (for $a=4$ about $64\%$ of
the mass is negative and $36\%$ positive, the integrand reaching $+2.78$ on the
positive part), and
membership of a contour cell in $\Omega_F=\{B_F>A_F\}$ is uncertain on straddle cells.
A sound upper bound on the (true negative) integral is accumulated cell-wise as
follows. On a cell, the outward-rounded enclosures $|R_a|^2\in[\rho^a_-,\rho^a_+]$,
$|R_{a+1}|^2\in[\rho^{a+1}_-,\rho^{a+1}_+]$ (the same $\rho$-machinery as
Section~\ref{sec:bound}) give half-log enclosures $\log|R_a|\in[\ell_a,L_a]$,
$\log|R_{a+1}|\in[\ell_{a+1},L_{a+1}]$, and the signed integrand is enclosed by the
directed-rounded difference $\bigl[\ell_{a+1}-L_a,\;L_{a+1}-\ell_a\bigr]$ (the
subtraction rounds each endpoint outward). Writing $h$ for the upper endpoint
$L_{a+1}-\ell_a$ and $w$ for the cell width, the cell's contribution to
$D_a^{\mathrm{hi}}$ is
\[
  \begin{cases}
    w\,h, & \text{if the cell is certainly in } \Omega_F\ (B_F\text{-lower}>A_F\text{-upper}),\\[2pt]
    w\,\max(0,h), & \text{if the cell straddles } \partial\Omega_F,\\[2pt]
    0, & \text{if the cell is certainly out } (B_F\text{-upper}<A_F\text{-lower}).
  \end{cases}
\]
The asymmetry is essential: on a certainly-in cell a negative $h$ correctly
\emph{reduces} the bound, but a straddle cell may lie outside $\Omega_F$, so its
negative integrand is discarded and only a positive part is banked. (Bounding the
integrand by $|\,\cdot\,|$ instead, as in the symmetric Lemma~\ref{lem:transfer}
accumulation, would round toward $0$ and is unsound for a negative integral.) The
per-cell upper endpoint $w\,h$ is loose only where $R_a,R_{a+1}$ share a deep
near-contour zero, but there $w\,h$ grows merely like $w\log(1/w)\to0$ under bisection;
adaptively bisecting every could-be-in cell whose contribution exceeds a fixed
tolerance (and every cell on which $|R_a|^2$ or $|R_{a+1}|^2$ is not yet certified
positive, securing $\inf_{\Omega_F}|R_a|>0$ for each $a$) drives $D_a^{\mathrm{hi}}$ to
the true difference. The $R_{a+1}/R_a$ ratio has milder wells than $R/Q_1$ (the two
siblings share most of their roots), so the harness of Lemma~\ref{lem:transfer} clears
it. Certificate \texttt{firstvar\_09} returns, at base grid $N=40000$ (reported in the
per-degree normalization $D_a/28$, whose sign is that of $D_a$),
\[
  D_1^{\mathrm{hi}}/28=-2.25\times10^{-3},\;
  D_2^{\mathrm{hi}}/28=-1.97\times10^{-3},\;
  D_3^{\mathrm{hi}}/28=-1.71\times10^{-3},\;
  D_4^{\mathrm{hi}}/28=-4.33\times10^{-4},
\]
all strictly negative, with $\inf_{\Omega_F}|R_a|>0$ certified for $a=1,\dots,5$ and no
unresolved cells. (The true per-degree differences are $-6.86,-5.25,-3.56,-1.84$ in
units of $10^{-3}$; the certified bounds over-estimate them, the safe direction for
proving $D_a<0$.) The four signs chain to the ordering.
\end{proof}

\begin{remark}[The unconstrained optimum sacrifices degree preservation]
\label{rem:a6}
The restriction to $a\le5$ is not cosmetic. For the subtraction family
$R_a=\pp(Q_1-g_a)$ the exponent $a=6$ gives a sibling whose per-degree active potential
$\rb{R_6}\approx-0.0523$ on $\Omega_F$ sits below $\rb{R_5}\approx-0.0498$, so ``$a=5$ is
firing-optimal'' is false without the restriction. But here
$\deg R_6=27$ (Proposition~\ref{prop:degfloor}): $a=6$ breaks degree preservation, the
shared leading coefficients of $Q_1$, and---decisively for the argument
above---the equal-degree hypothesis $\deg R_{a+1}=\deg R_a$ that makes the
shared-normalizer difference identity \eqref{eq:diff-id} hold, so no $D_5$ for the
$a=5\leftrightarrow a=6$ pair even exists. Proposition~\ref{prop:restricted-opt} is
therefore the sharpest defensible optimality statement, and it explains why Grinsztajn
stopped at $a=5$: $a=5$ is the firing-best choice consistent with degree preservation
and the transfer identity that certifies firing.
\end{remark}

%==============================================================================
\section{Computational evidence for saturation}\label{sec:saturation}
%==============================================================================

We finally use the base criterion of Corollary~\ref{cor:criterion} as a design objective and ask which additional integer blocks appear capable of improving the certified family. The blocks of Section~\ref{sec:bound} enter the base branch,
so the relevant quantity is the base marginal
$\widetilde r_Q=\int_{\{A_0>B\}}\log|Q(\chi)|\,ds$, with $Q$ firing iff
$\widetilde r_Q<0$. The results of this section are \emph{computational evidence} over
an explicit finite search, not a non-existence theorem; we are explicit about the
search's boundaries.

The one structural fact is additivity. Since
$\widetilde r_Q=\int_{\{A_0>B\}}\log|Q|\,ds$ is an integral of $\log|Q|=\sum\log|Q_i|$
over a fixed arc,
\begin{equation}\label{eq:additive}
  \widetilde r_{Q_1Q_2}=\widetilde r_{Q_1}+\widetilde r_{Q_2}.
\end{equation}
A reducible block fires only if it has a firing factor.

\begin{remark}[The cheap firing blocks are already in the dictionary]\label{rem:evidence}
Evaluated at the certified family of Section~\ref{sec:bound} on the candidate-free
anchor, an enumeration of Flammang's \citep{F18} Table~1 finds the firing entries
($\widetilde r_Q<0$) and sorts them by whether they are coprime to the active
dictionary $W=\{P_1,P_2,P_4,P_6,P_8,j_3,j_9,Q_1,Q_2,Q_5,Q_6\}$. The low-degree firing
entries that are \emph{not} coprime---degrees $1,3,12,12,16$---each share a factor with
$W$ and so are already absorbed; the four cheapest firing factors are exactly the base
blocks $j_3,j_9$ and the perturbing blocks $Q_5,Q_6$ used in the bound. Three
higher-degree firing table entries (degrees $16,16,20$, base marginals $\approx-0.005$
to $-0.008$) \emph{are} coprime to $W$: these would fire, but their realized gain after
re-optimization falls below the certificate's resolution (Section~\ref{sec:bound}), so
adjoining them does not yield a further certified break. A direct closest-vector and
root-target search over original integer blocks (degree $\le10$, bounded coefficients;
roughly $7.4\times10^4$ candidates, extended to degree $16$ near the firing locus)
returns no original block whose base marginal is more negative than these, and by
\eqref{eq:additive} a reducible block built from a coprime factor of cost
$\widetilde r\ge+0.013$ (the cheapest such factor found, $X^2-X+1$) cannot beat the
best firing margin $\approx-0.005$. The evidence is therefore that the cheap,
\emph{certifiably} improving blocks of this construction have been absorbed---not that
no admissible block of any degree fires.
\end{remark}

\begin{remark}[Why the cheap firing set is so thin]\label{rem:wells}
The firing integer blocks are isolated in the coefficient lattice: a single $\pm1$
change to a firing table entry swings its base marginal from $\approx-0.005$ to above
$+0.06$. Geometrically, $\widetilde r_Q$ is negative only when a root of $Q$ sits in a
deep well of the base active-arc potential, and those wells lie strictly inside the
lemniscate $\chi(\{|z|<1\})$, where an admissible integer block cannot place a root
without a conjugate elsewhere that pays for it. This is the mechanism behind the
thinness; the statements above rest on the additivity \eqref{eq:additive} and the
reproduced marginals, not on the geometry.
\end{remark}

\begin{remark}[Not a transfinite-diameter problem]\label{rem:noteq}
Factoring $\log|Q|$ over the roots writes $\rQ$ as an averaged logarithmic potential
of $\Omega$, which invites reading block selection as an equilibrium
(transfinite-diameter) problem. It is not: the equilibrium measure of $\Omega$ has
constant potential $\approx0.524$ on its support, well above $\log h=0.254$, so
$\Omega$'s measure is not the equilibrium measure and $\log h$ is not its Robin
constant. The integer constraint, not capacity, controls the value, and the correct
frame is the first-variation calculus above.
\end{remark}

%==============================================================================
\section{Certification and reproducibility}\label{sec:verify}
%==============================================================================

Every numerical assertion above is reproduced by a self-contained program in the
repository \repo, in the directory \texttt{constants/82a/certificate/}. The claims fall
into two kinds, and we are explicit about which is which. The \emph{certified} claims---
the improved bound (Theorem~\ref{thm:upper}), the sibling-generator firing
(Theorem~\ref{thm:generator}, Lemma~\ref{lem:transfer}), and the bridge-exponent ordering
(Proposition~\ref{prop:restricted-opt})---are established by directed (outward) interval
arithmetic with an explicit accounting of floating-point round-off in the cell sums
(below), and each rejects a deliberately understated target with the frontier resolved,
ruling out a grid fallback. The \emph{grid diagnostics}---the record marginals of
Proposition~\ref{thm:record-fires}---are fine-grid evaluations, not interval certificates,
and are labelled as such in their statement; the programs reproduce their displayed
values but do not enclose them rigorously. The scientific Python stack
(\texttt{numpy}, \texttt{scipy}, \texttt{mpmath}, \texttt{sympy}) suffices.
Table~\ref{tab:verify} keys each statement to its check, marking the grid diagnostics.

\begin{table}[ht]
\centering\small
\renewcommand{\arraystretch}{1.15}
\begin{tabular}{@{}lll@{}}
\toprule
statement & program & verdict\\
\midrule
Thm.~\ref{thm:firstvar} & \texttt{firstvar\_01\_lemma} & sign $=$ fired/dry on all blocks\\
\quad hypotheses (L), (N) & \texttt{firstvar\_02\_hypotheses} & $K$ finite, dominator $L^1$\\
Thm.~\ref{thm:unified} & \texttt{firstvar\_03\_unified} & three regimes \texttt{PASS}, tie kink\\
Cor.~\ref{cor:criterion} & \texttt{firstvar\_04\_perturbing\_marginal} & marginal matches finite difference\\
Prop.~\ref{thm:record-fires} & \texttt{firstvar\_07\_record\_blocks} & fires on all 6, dry on controls\\
Thm.~\ref{thm:generator}, Lem.~\ref{lem:transfer} & \texttt{firstvar\_08\_sibling\_generator} & firing-transfer bound clears gap\\
\quad degree floor (Prop.~\ref{prop:degfloor}) & \texttt{firstvar\_08\_sibling\_generator} & $\deg R_a{=}28\!\iff\!a{\le}5$, admissible\\
Prop.~\ref{prop:restricted-opt} & \texttt{firstvar\_09\_restricted\_optimality} & $D_a^{\mathrm{hi}}{<}0$, $a{=}1..4$: $\rb{R_a}$ $\downarrow$, opt.\ $a{=}5$\\
\quad dictionary admissibility & \texttt{firstvar\_06\_dictionary} & pairwise coprime, squarefree\\
\S\ref{sec:saturation} (evidence) & \texttt{firstvar\_05\_numerator\_split} & active-set decomposition\\
\quad finite firing search & (saturation search) & cheapest new factor $\widetilde r=+0.013$\\
Thm.~\ref{thm:upper} & \texttt{bound\_07\_block\_j9} & \texttt{CERTIFIED} $\le0.2538893183$\\
\quad anchor chain & \texttt{bound\_00}--\texttt{bound\_06} & each $q_k{=}0$ recovers the previous\\
\bottomrule
\end{tabular}
\medskip
\caption{The verification set. The bound chain \texttt{bound\_00}--\texttt{bound\_07}
builds the enriched family one block at a time, each certificate reducing to the
previous when its new exponent is set to zero; the \texttt{firstvar} programs check
the first-variation theory and the selection criterion. The saturation enumeration and
the potential-theoretic computations are in the \texttt{scratch/} subdirectory.}
\label{tab:verify}
\end{table}

The certified bound is produced by \texttt{bound\_07\_block\_j9} in \texttt{certify}
mode on the exponent vector of Section~\ref{sec:bound}; the full command and the
chain of anchor checks are documented in the directory's \texttt{README}. The
results of this paper correspond to the repository state at the tagged release commit
recorded with the archived artifact accompanying the submission; the
polynomial data on which Theorem~\ref{thm:upper} depends is reproduced in full in
Appendix~\ref{app:data} so that the certificate does not rely on a mutable link.

%==============================================================================
\appendix
\section{Polynomial data}\label{app:data}
%==============================================================================

All polynomials are written in the symmetric coordinate $X=z(1-z)$, by descending
coefficients. The base family and distinguished block are Doche's \citep{Doc01b}; the
low-degree base polynomials are
\[
  P_1=X,\qquad P_2=1-X,\qquad P_4=X^4-2X^3+4X^2-3X+1,
\]
and the higher base polynomials, together with the two degree-$28$ factors $Q_1,Q_2$
of the distinguished block, are listed by coefficient vector:
{\small
\begin{align*}
P_6:&\ [1,-3,8,-16,26,-27,17,-6,1],\\
P_8:&\ [1,-3,7,-14,30,-58,96,-123,114,-73,31,-8,1],\\
Q_1:&\ [1,-7,30,-97,269,-679,1612,-3618,7646,-15180,28457,-50741,86189,\\
   &\ \ -138288,206152,-279897,339335,-360911,331775,-260367,172556,\\
   &\ \ -95554,43677,-16221,4786,-1084,178,-19,1],\\
Q_2:&\ [1,-7,30,-96,255,-586,1212,-2360,4573,-9148,18749,-37783,71770,\\
   &\ \ -124910,195848,-273368,335981,-359545,331349,-260271,172542,\\
   &\ \ -95553,43677,-16221,4786,-1084,178,-19,1].
\end{align*}}
The perturbing blocks $Q_5,Q_6$ and the two base blocks $j_3,j_9$ adjoined in
Section~\ref{sec:bound} are Flammang \citep{F18} Table~1 entries:
{\small
\begin{align*}
Q_5=j_{13}:&\ [1,-3,8,-18,36,-62,97,-123,114,-73,31,-8,1]\quad(\deg 12),\\
Q_6=j_{15}:&\ [1,-4,10,-17,26,-47,119,-298,592,-878,963,-780,464,-199,59,-11,1]\quad(\deg 16),\\
j_3:&\ [1,1,-2,1]\quad(\deg 3),\qquad
j_9:\ [1,-1,0,-3,15,-22,16,-6,1]\quad(\deg 8).
\end{align*}}
The two near-cancellation factors of the \citet{Gri26} record
(Section~\ref{sec:record}), defined by \eqref{eq:R0R2} and taken to a positive
leading coefficient, are the degree-$28$ polynomials
{\small
\begin{align*}
R_0:&\ [1,-7,31,-107,321,-867,2148,-4922,10497,-20908,38973,-68026,111075,\\
   &\ \ -169111,238522,-308360,360033,-373190,337618,-262546,173171,-95678,\\
   &\ \ 43693,-16222,4786,-1084,178,-19,1],\\
R_2:&\ [1,-7,29,-86,202,-387,618,-861,1256,-2589,7105,-19328,45963,-93546,\\
   &\ \ 163248,-244838,315271,-347265,325506,-258092,171927,-95429,43661,\\
   &\ \ -16220,4786,-1084,178,-19,1].
\end{align*}}
Here $P_7,P_9$ in the record's denominator are the polynomials denoted $Q_5,Q_6$
above.

%==============================================================================
\bibliographystyle{plainnat}
\bibliography{references}

\end{document}
